\begin{titlepage}
\centering
{\Large Notas didácticas de una entrevista a fondo\par}
\vspace{1.0cm}
{\huge\bfseries Apostar por OpenAI, Robinhood y las cuentas del capital en IA\par}
\vspace{0.5cm}
{\Large Zhang Xiaojun conversa con Freda, socia de Altimeter\par}
\vspace{0.35cm}
{\large Elaborado a partir del video público de Zhang Xiaojun Podcast, una transcripción hecha con GPU, la estructura de capítulos y diagramas conceptuales propios\par}
\vspace{0.3cm}
{\large 2025-12-16\par}
\vspace{0.85cm}
\includegraphics[width=0.88\textwidth,height=0.42\textheight,keepaspectratio]{cover.jpg}\par
\vfill
\begin{tcolorbox}[width=0.92\textwidth, colback=black!2!white, colframe=black!55, sharp corners]
\textbf{Título del video}: 125. Conversación con Freda, socia de Altimeter: la apuesta por OpenAI, la historia de Robinhood, los «chicos malos» del capital estadounidense, los cálculos y la burbuja\par
\textbf{Canal del video}: Zhang Xiaojun Podcast\par
\textbf{Duración del video}: 01:24:26\par
\textbf{Enlace del video}: \href{\videourl}{\nolinkurl{\videourl}}\par
\textbf{Nota sobre los materiales}: YouTube no ofrece subtítulos disponibles; el texto se elaboró a partir de una transcripción con faster-whisper large-v3 en CUDA, los capítulos del video, la portada y figuras didácticas propias. Las tomas fijas de la entrevista no se reutilizan en el texto.\par
\textbf{Página del proyecto}: \href{\repourl}{\nolinkurl{\repourl}}\par
\end{tcolorbox}
\end{titlepage}

\tableofcontents
\newpage
